% ============================================================================
% Computational and structural bounds for the Erdos--Gyarfas conjecture
% Archived draft mathematical note. Compile: pdflatex x2.
%
% Archival bibliographic limitations: the prize/status and independent
% forum-observation date were recorded in campaign notes, not independently
% checked against the live site. The P8/P10 references are abbreviated.
% Computational status is frozen through the independently verified n=48 rung.
% ============================================================================
\documentclass[11pt]{article}

\usepackage[margin=1.1in]{geometry}
\usepackage{amsmath,amssymb,amsthm}
\usepackage{booktabs}
\usepackage[colorlinks=true,linkcolor=blue,citecolor=blue,urlcolor=blue]{hyperref}

\newtheorem{theorem}{Theorem}[section]
\newtheorem{lemma}[theorem]{Lemma}
\newtheorem{corollary}[theorem]{Corollary}
\newtheorem{conjecture}[theorem]{Conjecture}
\theoremstyle{remark}
\newtheorem{remark}[theorem]{Remark}

\newcommand{\EG}{Erd\H{o}s--Gy\'arf\'as}
\newcommand{\repourl}{\url{https://github.com/mcolsen/erdos-gyarfas}}
\newcommand{\Cyc}[1]{C_{#1}}

\title{Computational and structural bounds for the\\ \EG{} conjecture}
\author{Maria Olsen}
\date{August 8, 2026}

\begin{document}
\maketitle

\begin{abstract}
The \EG{} conjecture asserts that every graph with minimum degree at least~$3$
contains a cycle whose length is a power of two. We report a computational
campaign against the conjecture, using independent software stacks or runs
where feasible and count-exact nonzero controls otherwise. Main
results: (i)~by isomorph-free exhaustive generation, \emph{every} graph with
$\delta \ge 3$ on $n \le 21$ vertices contains a cycle of length $4$, $8$,
or~$16$ (the $n=21$ order was enumerated twice, via two complete
res/mod partitions, $\approx 900$ CPU-hours total); (ii)~by SAT with
isomorph-free model enumeration, there is no counterexample on $n \le 32$
vertices, and, adding sound structural constraints on a minimum-order
counterexample, any counterexample requires \textbf{at least $36$ vertices};
(iii)~any \emph{cubic} counterexample requires at least $50$ vertices, and any
\emph{bipartite} counterexample at least $57$; (iv)~at least $2/3$ of the
vertices of a minimal counterexample have degree exactly~$3$, strengthening
the recent $4/7$ bound of Carr, and every $\Cyc4$-free graph with
$\delta \ge 3$ satisfies $\Delta \le \lfloor (n-1)/2 \rfloor$; (v)~all
$3{,}815$ cubic arc-transitive graphs on at most $10{,}000$ vertices and all
$111{,}360$ cubic vertex-transitive graphs on at most $1{,}280$ vertices
satisfy the conjecture; (vi)~the smallest cubic graph avoiding
$\{\Cyc4,\Cyc8,\Cyc{16}\}$ has between $50$ and $82$ vertices; (vii)~a
Moore-type bound: a cubic graph whose non-backtracking transfer matrix
satisfies $\operatorname{tr}(H^L)=0$ has at least $\approx 2^{L/2-1}$
vertices, which rules out entire families of algebraic constructions at small
order. We close with a first-moment analysis explaining \emph{why} the
conjecture resists counterexamples --- the expected number of $\Cyc L$'s in a
cubic-like graph grows doubly exponentially in the index of $L$ --- and a
precise characterization of the only kind of object that could still refute
it. Source, compact execution records, and graph artifacts are archived at
\repourl.
\end{abstract}

\section{Introduction}

Erd\H{o}s and Gy\'arf\'as conjectured in the mid-1990s that every graph of
minimum degree at least three contains a simple cycle whose length is a power
of two~\cite{EG95}. Erd\H{o}s offered \$100 for a proof and \$50 for a
counterexample, and reportedly considered the conjecture likely to be
\emph{false}; it is problem~\#64 on \texttt{erdosproblems.com}, where a
\$1000 prize was recorded in the campaign notes.
% The prize/status line was not independently checked against the live site.

The conjecture is known to hold for several restricted classes: planar
claw-free graphs (Daniel--Shauger~\cite{DS01}), $3$-connected cubic planar
graphs (Heckman--Krakovski~\cite{HK13}), $K_{1,m}$-free graphs with minimum
degree $m{+}1$ or maximum degree $2m{-}1$~\cite{Shauger98}, and graphs with
no long induced path ($P_8$-free~\cite{P8free}, $P_{10}$-free~\cite{P10free}).
In a different direction, the cycle-length machinery of Liu and
Montgomery~\cite{LM23} implies that graphs of sufficiently large average
degree contain cycles of every even length in a long interval, hence in
particular a power-of-two cycle; the conjecture is thus ``only'' open in the
sparse regime it was designed for. On the computational side the record
stood, until recently, at exhaustive verification for all graphs on roughly
$15$--$16$ vertices (attributed to Royle, c.~2004) and for cubic graphs on
fewer than $30$ vertices (Markstr\"om~\cite{Markstrom04}). In 2026 two
unrefereed contributions appeared: an SMS-based SAT frontier claiming no
counterexample on $n \le 31$ vertices (Balaji~\cite{Balaji26}), and a short
structural paper of Carr~\cite{Carr26} showing every minimal counterexample
is ``predominantly cubic.'' Part of the present campaign was an independent
adjudication of both; both are sound (Section~\ref{sec:sat}), and Carr's
bound can be strengthened (Section~\ref{sec:structural}).

\subsection{Results}

Throughout, a \emph{counterexample} is a graph $G$ with $\delta(G)\ge 3$
containing no cycle whose length is a power of two (equivalently, of length
$2^k$ for some $k \ge 2$). A \emph{minimal}
counterexample is one of minimum order, and of minimum size among those of
minimum order.

\begin{enumerate}
\item \textbf{Exhaustive frontier (no SAT in the loop).} By isomorph-free
exhaustive generation (\texttt{geng} with a custom pruning plugin), there are
exactly $0$ graphs with $\delta \ge 3$ avoiding $\{\Cyc4, \Cyc8, \Cyc{16}\}$
at \emph{every} order $n = 4, \dots, 21$. Since $32 > 21$, this verifies the
conjecture on $n \le 21$ outright. The last order was enumerated twice, by
two complete and independently sliced work partitions in exact agreement
(Section~\ref{sec:enum}).
\item \textbf{SAT frontier.} Independently of (1), SAT-with-symmetries
enumeration proves there is no counterexample on $n \le 32$ vertices,
assumption-free. On top of this base, sound minimum-order structural
constraints (Section~\ref{sec:structural}) extend the bound: \emph{any
counterexample has at least $36$ vertices}. Any cubic counterexample has at
least $50$ vertices; any bipartite counterexample at least $57$
(Section~\ref{sec:sat}).
\item \textbf{Structural bounds.} At least $2/3$ of the vertices of a
minimal counterexample are cubic (Theorem~\ref{thm:twothirds}, strengthening
Carr's $4/7$); and every $\Cyc4$-free graph with $\delta\ge3$ has
$\Delta \le \lfloor (n-1)/2\rfloor$ (Theorem~\ref{thm:degcap} --- no
minimality needed, so it prunes \emph{every} counterexample search).
\item \textbf{Symmetric censuses are clean.} All $3{,}815$ cubic
arc-transitive graphs on $\le 10{,}000$ vertices (Conder census) and all
$111{,}360$ cubic vertex-transitive graphs on $\le 1{,}280$ vertices
(Poto\v{c}nik--Spiga--Verret census) satisfy the conjecture. The interesting
part is \emph{how}: $689$ resp.\ $100$ of them avoid
$\{\Cyc4,\Cyc8,\Cyc{16}\}$ entirely --- yet every single one contains
$\Cyc{32}$ (Section~\ref{sec:census}).
\item \textbf{A new extremal question, bracketed.} The smallest cubic graph
with no $\Cyc4$, $\Cyc8$, or $\Cyc{16}$ has between $50$ and $82$ vertices
(Section~\ref{sec:window}); eleven verified specimens at $n = 82$--$126$ are
archived.
\item \textbf{A Moore-type obstruction for algebraic constructions.} If a
connected cubic graph has $\operatorname{tr}(H^L) = 0$ for the Hashimoto
(non-backtracking) matrix $H$ and even $L$, then
$n \ge (s+1)^2/(2s-1)$ with $s = 2^{L/2}$ (Section~\ref{sec:outlook});
certifying $\Cyc{64}$-freeness by walk cancellation therefore needs
$n > 2.1\times 10^9$. This kills, at realistic orders, the most natural
family of algebraic constructions (voltage lifts certified through twisted
spectra).
\item \textbf{A margin heuristic and a target profile.} The expected number
of $L$-cycles in a cubic-like graph is $\approx 2^L/(2L)$ --- doubly
exponential in $k$ where $L = 2^k$. Everywhere we could measure, observed
behavior tracked this formula, and \emph{spectral holes occurred only at
lengths where the expected count is small}. Any counterexample must be a
highly structured algebraic object with engineered holes at two or more
consecutive large powers of two against millions of expected cycles ---
an object nobody currently knows how to build (Section~\ref{sec:outlook}).
\end{enumerate}

\subsection{Methodology: adversarial cross-validation}
\label{sec:method}

Computer-search claims of the form ``zero objects exist'' are only as strong
as their weakest unverified component. The campaign's standing rule was that
every headline zero must be reproduced by two \emph{disjoint} stacks --- or,
where cost forbade that, that every ingredient of the single stack be
validated count-exactly against an independent implementation on nonzero
instances. Concretely: the \texttt{nauty}/\texttt{geng} stack (C, canonical
augmentation) and the SMS stack (SAT + minimal-model symmetry breaking +
subgraph-isomorphism propagator) were run against each other on five
count-exact gates (agreeing at $2590$, $5$, $1315$, $0$, and $88$ models;
Section~\ref{sec:sat}); the exact cycle checker was validated against an
algorithmically independent brute force on all $13{,}591$ graphs with
$n \le 8$ and against published spectra of named graphs; and the enumeration
reproduced Markstr\"om's published counts ($4$ and $23$ cubic
$\{\Cyc4,\Cyc8\}$-free graphs at $n = 24, 26$) before being trusted at new
orders. One coverage error in an early census sweep was caught by exactly
this discipline and is documented in the repository (the corrected census
results are those stated in Section~\ref{sec:census}).

\section{Structural results}
\label{sec:structural}

Following Carr~\cite{Carr26}, whose short proofs we verified independently
and reuse, let $G$ be a minimal counterexample (minimum order, then minimum
size). Write $V_3$ for its set of degree-$3$ vertices and $V_{\ge4}$ for the
rest.

\begin{lemma}[{Carr~\cite{Carr26}}]\label{lem:proper}
Every proper subgraph $H \subsetneq G$ has $\delta(H) \le 2$.
\end{lemma}

\begin{proof}
$H$ precedes $G$ in the (order, size) lexicographic order; if
$\delta(H)\ge3$, minimality forces a power-of-two cycle in $H \subseteq G$.
\end{proof}

\begin{corollary}[{Carr~\cite{Carr26}}]\label{cor:cubicneighbor}
Every vertex of $G$ is adjacent to a vertex of degree exactly $3$.
\end{corollary}

\begin{proof}
$G - v$ is proper, so some $u$ has $\deg_{G-v}(u)\le 2$; since
$\deg_G(u)\ge3$, $u$ lost an edge to $v$ and $\deg_G(u)=3$.
\end{proof}

\begin{corollary}[{Markstr\"om~\cite{Markstrom04}}]\label{cor:indep}
$V_{\ge4}$ is an independent set.
\end{corollary}

\begin{proof}
Deleting an edge between two vertices of degree $\ge4$ leaves a proper
subgraph with $\delta\ge3$, contradicting Lemma~\ref{lem:proper}.
\end{proof}

Carr combines these into: at least $4/7$ of the vertices are cubic. His own
two corollaries give more in one line.

\begin{theorem}\label{thm:twothirds}
In a minimal counterexample, $|V_3| \ge \tfrac{2}{3}\,|V(G)|$.
\end{theorem}

\begin{proof}
By Corollary~\ref{cor:cubicneighbor} applied to the members of $V_3$, every
cubic vertex has at least one cubic neighbor, so each $v \in V_3$ sends at
most $2$ of its $3$ edges to $V_{\ge4}$. By Corollary~\ref{cor:indep}, all
edges at $V_{\ge4}$ land in $V_3$. Counting the edges between the two sides:
\[
4\,|V_{\ge4}| \;\le\; e(V_3, V_{\ge4}) \;\le\; 2\,|V_3|
\quad\Longrightarrow\quad
|V_{\ge4}| \le \tfrac{1}{2}|V_3|
\quad\Longrightarrow\quad
|V_3| \ge \tfrac{2}{3}|V(G)|. \qedhere
\]
\end{proof}

\begin{remark}
Equality would force every $V_{\ge4}$-vertex to have degree exactly $4$,
and $V_3$ to induce a perfect matching with each cubic vertex sending
exactly two edges to $V_{\ge4}$.
The same bound was observed independently, and essentially simultaneously,
in a forum post by user \emph{jul059} on \texttt{erdosproblems.com}
(26~July 2026); the proof above was recorded in this project's repository on
24~July 2026. % Forum attribution/date retained from the campaign record.
\end{remark}

The second structural bound requires no minimality at all --- only
$\Cyc4$-freeness, which every counterexample has since $4 = 2^2$. It is what
makes exhaustive generation feasible at orders around $20$
(Section~\ref{sec:enum}).

\begin{theorem}\label{thm:degcap}
If $G$ is $\Cyc4$-free with $\delta(G) \ge 3$ and $n$ vertices, then for
every vertex $v$,
\[
\sum_{u \in N(v)} \bigl(\deg(u) - 1\bigr) \;\le\; n - 1,
\]
and consequently $\Delta(G) \le \lfloor (n-1)/2 \rfloor$.
\end{theorem}

\begin{proof}
Fix $v$. The sets $N(u)\setminus\{v\}$ for $u \in N(v)$ are pairwise
disjoint: a common element $w$ of two of them would close a $4$-cycle
$v u_1 w u_2$. They avoid nothing else, so their total size,
$\sum_{u\in N(v)}(\deg(u)-1)$, is at most $|V(G)\setminus\{v\}| = n-1$.
For the consequence, take $v$ of maximum degree $\Delta$: since every
$u \in N(v)$ has $\deg(u) \ge 3$, each of the $\Delta$ terms is at least
$2$, so $2\Delta \le n-1$.
\end{proof}

\section{Exhaustive generation: no counterexample on $n \le 21$ vertices}
\label{sec:enum}

The generator is McKay's \texttt{geng}~\cite{nauty} with minimum degree $3$,
maximum degree $\lfloor (n-1)/2 \rfloor$ (sound by
Theorem~\ref{thm:degcap}), the built-in square-free filter (\texttt{-f},
sound since $\Cyc4$ is itself a forbidden cycle), and a custom
\texttt{PRUNE} plugin rejecting any partial graph containing a $\Cyc8$ or
$\Cyc{16}$ through the newest vertex. The plugin is inductively exact: in
canonical augmentation each graph is built by successive vertex additions,
and any forbidden cycle in the final graph passes through the
highest-indexed vertex it contains at the moment that vertex was added ---
so it is detected at that step, and each isomorph-free branch is pruned at
the earliest opportunity and never wrongly. The plugin reproduced
Markstr\"om's published Table (the $4$ and $23$ cubic $\{\Cyc4,\Cyc8\}$-free
graphs at $n=24,26$, each of which we verified contains a $\Cyc{16}$) before
any new claim was made with it.

\textbf{Result.} There are exactly $0$ graphs with $\delta\ge3$ avoiding
$\{\Cyc4,\Cyc8,\Cyc{16}\}$ at every order $n = 4,\dots,21$. Every graph with
minimum degree at least $3$ on at most $21$ vertices therefore contains a
cycle of length $4$, $8$, or $16$. This extends the historical
pure-enumeration record ($\approx 15$--$16$ vertices) by five to six orders,
with no SAT solver in the loop.

The final order $n=21$ was distributed with \texttt{geng}'s res/mod
splitting and completed \emph{twice}: once as a complete $256$-class
partition ($509$ CPU-h) and once as a complete $16$-class partition ($389$
CPU-h), sliced along different tree boundaries, run on different days, in
exact agreement (zero survivors in every class of both partitions).

\begin{remark}[A soundness warning on res/mod splitting]\label{rem:resmod}
It is tempting to refine a slow class $i \pmod m$ into the classes
$i + km \pmod{4m}$ of a finer run. This is \emph{unsound} for
\texttt{geng}: the level at which the generation tree is split is adjusted
dynamically by a threshold proportional to \texttt{mod} itself
(\texttt{multiplicity} $= 50\cdot$\texttt{mod} in \texttt{geng.c}), so runs
with different moduli can split at different tree levels, and classes of
different moduli need not be compatible --- the union above can both miss
and double-count work. Any \emph{single} modulus, run over all its residues
with identical options, partitions the tree exactly (all such runs share
identical traversal state above the split level). Our two certifications
are each pure single-modulus partitions; we observed empirically that
matching residue classes across moduli mis-align badly (per-class workloads
differing by factors of $0.16$ to $2.2$), so the warning is not
theoretical.
\end{remark}

\section{SAT frontiers: general $\ge 36$, cubic $\ge 50$, bipartite $\ge 57$}
\label{sec:sat}

The SAT stack is SAT Modulo Symmetries (SMS) of Kirchweger and
Szeider~\cite{SMS}, which enumerates isomorph-free models under a
lex-minimality symmetry break, with the Glasgow subgraph
solver~\cite{Glasgow} as a forbidden-subgraph propagator (forbidding the
cycles $\Cyc4$, $\Cyc8$, $\Cyc{16}$, and where relevant $\Cyc{32}$, as
non-induced subgraphs). We first audited the 2026 frontier claim of
Balaji~\cite{Balaji26} ($n \le 31$): the encoding-level design is sound
(blocking clauses, symmetry breaking, and relaxation direction all verified
by code review), its published test gates pass, and we replicated the entire
frontier $n = 17$--$31$ (all UNSAT) on our hardware. Since the stack's
correctness cannot rest on proof logging (the propagator's clauses are not
RUP-derivable), we validated it count-exactly against the independent
\texttt{geng} stack of Section~\ref{sec:enum} on five gates:

\begin{center}
\begin{tabular}{lccc}
\toprule
instance & \texttt{geng} count & SMS count & agree \\
\midrule
$n=8$, $\delta\ge3$, nothing forbidden & 2590 & 2590 & \checkmark \\
$n=10$, forbid $\Cyc4$ & 5 & 5 & \checkmark \\
$n=10$, forbid $\Cyc8$ & 1315 & 1315 & \checkmark \\
$n=12$, forbid $\{\Cyc4,\Cyc8\}$ & 0 & 0 & \checkmark \\
$n=16$, forbid $\{\Cyc4,\Cyc{16}\}$ & 88 & 88 & \checkmark \\
\bottomrule
\end{tabular}
\end{center}

\textbf{Assumption-free base.} SMS proves there are zero
$\{\Cyc4,\Cyc8,\Cyc{16}\}$-free graphs with $\delta \ge 3$ at $n = 32$
(5.5\,h) --- the first order at which a $\Cyc{32}$ could exist, and it is
not even needed. Combined with $n \le 31$: \emph{no counterexample exists on
at most $32$ vertices, with no structural assumptions.}

\textbf{Structural extension to $\ge 36$.} For $n = 33, 34, 35$ we add the
constraints of Section~\ref{sec:structural}, all sound for a
\emph{minimum-order} counterexample: $V_{\ge4}$ independent
(Corollary~\ref{cor:indep}), every vertex adjacent to a cubic vertex
(Corollary~\ref{cor:cubicneighbor}), $|V_{\ge4}| \le n/3$
(Theorem~\ref{thm:twothirds}), and the degree cap
(Theorem~\ref{thm:degcap}). The chaining argument: if any counterexample
existed on $\le 35$ vertices, the global minimum-order one would live at
some order $\le 35$ and satisfy all constraints; orders $\le 32$ are
excluded unconditionally, and the constrained searches at $33, 34, 35$
(5\,h, 10.5\,h, 22\,h) are UNSAT. Hence \textbf{any counterexample has at
least $36$ vertices}. (The encoding was validated by re-running orders with
known nonzero unconstrained counts and checking the constrained counts
against \texttt{geng} filters.)

\textbf{Cubic.} For cubic graphs the SMS ladder is assumption-free at each
even order, and gives zero $\{\Cyc4,\Cyc8,\Cyc{16}\}$-free cubic graphs at
$n = 32, 34, \dots, 48$ (5m52s at $32$ growing to a 146-hour monolith at
$48$). The n=48 zero was separately reproduced by a complete recursive
cube partition after the cubing pipeline matched Markstr\"om's positive
n=28 count exactly at three granularities. With Markstr\"om's $n < 30$,
parity, and the general result $\le 32$: \textbf{any cubic counterexample has
at least $50$ vertices}.

\textbf{Bipartite.} Adding a $2$-coloring to the encoding (sound within the
class), zero bipartite $\{\Cyc4,\Cyc8,\Cyc{16}\}$-free graphs with
$\delta\ge3$ exist at every $n = 33, \dots, 56$: \textbf{any bipartite
counterexample has at least $57$ vertices}.

\section{Symmetric censuses and structured families}
\label{sec:census}

Exhaustion is hopeless much past $n \approx 45$, but a counterexample, if
one exists, is plausibly a highly symmetric object (Section~\ref{sec:outlook}
makes this precise). We therefore swept every available census and several
parametric families with the exact cycle checkers:

\begin{center}
\begin{tabular}{lrrl}
\toprule
family & tested & $\{\Cyc4,\Cyc8,\Cyc{16}\}$-free & fate of survivors \\
\midrule
Generalized Petersen $GP(n,k)$, $2n\le64$ & 240 & 0 & --- \\
$I$-graphs $I(n,j,k)$, $2n\le64$ & 1{,}148 & 0 & --- \\
Circulants, $n\le64$, degree 4--6 & 80{,}780 & 0 & --- \\
Cyclic lifts of $\theta, K_4, K_{3,3}, Q_3$, Petersen; $n\le64$ & 176{,}348 & 0 & --- \\
Conder arc-transitive census, $n\le10{,}000$ & 3{,}815 & 689 & all contain $\Cyc{32}$, $\Cyc{64}$, $\Cyc{128}$ \\
PSV vertex-transitive census, $n\le1{,}280$ & 111{,}360 & 100 & all contain $\Cyc{32}$, $\Cyc{64}$ \\
\bottomrule
\end{tabular}
\end{center}

Both censuses~\cite{Conder,PSV13} are \emph{clean}: the conjecture holds for
every cubic arc-transitive graph on at most $10{,}000$ vertices and every
cubic vertex-transitive graph on at most $1{,}280$. ($\Cyc{16}$-freeness
decisions were triple-engined --- exact DFS, non-backtracking dart search,
and Glasgow --- with zero disagreements; every claimed long cycle has an
explicitly verified witness.)

The survivors of the $\{\Cyc4,\Cyc8,\Cyc{16}\}$ pass are extremal objects of
independent interest. Three deserve mention:
\begin{itemize}
\item \textbf{C2030.1} ($n=2030$, girth $7$, non-bipartite; the skeleton of
a non-orientable $\{7,3\}$ map of $\mathrm{PSL}(2,29)$): it contains
$7$-cycles yet no $\Cyc8$ or $\Cyc{16}$ --- spectral holes without high
girth, at lengths where a random cubic graph would carry ${\sim}16$ and
${\sim}2000$ cycles. Notably its non-backtracking walk count at length
$16$ is $194{,}880 \ne 0$: the hole exists at the \emph{cycle} level while
being invisible at the \emph{walk} level.
\item \textbf{C1458.10/11} ($n=1458$, girth $12$, bipartite): even-cycle
spectra skipping $16$, and here the walk count vanishes too --- a genuinely
representation-theoretic cancellation.
\item \textbf{The smallest survivors}: C1050.1 ($n=1050$, girth $12$) in
the arc-transitive census, and an $n=180$, girth-$3$ (!) graph in the
vertex-transitive census --- sixty triangles and then no cycle of any
length in $\{4,\dots,11\}$.
\end{itemize}
Every one of the census survivors contains a $\Cyc{32}$, typically
found by the Glasgow solver in tens of milliseconds: the symmetry that
engineers the short holes \emph{spreads} the long cycles.

\section{Stochastic search, an extremal bracket, and the $[65,127]$ window}
\label{sec:window}

\textbf{Simulated annealing.} Over cubic graphs with energy
$10\,\#\Cyc4 + 5\,\#\Cyc8 + \#\Cyc{16}$ (exact counts per move; an exact
$\Cyc{32}$ gate on any zero-energy state), no chain among $30+$ across three
move strategies and ${\sim}1.5$ billion total moves ever reached energy $0$
at any $n \in [34,62]$: best energies floor at $13$--$25$, essentially
independent of $n$ and strategy. The conjecture's difficulty is not an
artifact of exhaustive methods; even aggressive local search cannot
approach a counterexample in this range.

\textbf{The bracket.} At $n \ge 82$ the picture inverts:
$\{\Cyc4,\Cyc8,\Cyc{16}\}$-free cubic graphs are \emph{easy} to find
(annealing succeeds in minutes at $n = 82, 90, 94, 96, 112, 120, 126$;
eleven specimens re-verified by an independent checker are archived). With
the SMS ladder zero through $48$: \emph{the smallest cubic graph avoiding
$\{\Cyc4, \Cyc8, \Cyc{16}\}$ has between $50$ and $82$ vertices.} Repeated
annealing probes at $n = 66$--$80$ all failed while succeeding instantly at
$82$, hinting the true threshold sits near the top of the bracket.

\textbf{The $[65,127]$ window.} In this window a counterexample must avoid
five lengths $\{4,8,16,32,64\}$. The specimens above show the first three
exclusions are cheap here --- but every specimen carries at least $10^5$
thirty-two-cycles, and girth $>32$ requires $n \gtrsim 1.3\times10^5$ by the
Moore bound, far beyond the window. A counterexample here would need
engineered spectral holes at $32$ \emph{and} $64$ simultaneously against
$10^5$-plus expected cycles of each length --- outside the reach of any
local or exhaustive methods used in this campaign. This is a heuristic
assessment, not an exhaustion of the window.

\section{Margin analysis, an algebraic obstruction, and limitations}
\label{sec:outlook}

\textbf{Why the conjecture keeps winning.} The expected number of $L$-cycles
in a random cubic-like graph is ${\sim}2^L/(2L)$: avoiding $\Cyc8$ means
dodging ${\sim}16$ expected cycles; $\Cyc{16}$, ${\sim}2\times10^3$;
$\Cyc{32}$, ${\sim}7\times10^7$; $\Cyc{64}$, ${\sim}10^{17}$. As $n$ grows
the forbidden lengths thin out (the reason Erd\H{o}s doubted the
conjecture), but each forbidden length gets \emph{doubly exponentially}
fatter --- and in every experiment above, the second effect won: the
$\{\Cyc4,\Cyc8\}$-free cubic counts explode ($4 \to 23 \to 251$ at
$n=24,26,28$) while adding the $\Cyc{16}$ exclusion flattens them to zero
through $48$; annealing floors at a persistent residue of $\Cyc{16}$'s; in
the window, the wall simply moves to $\Cyc{32}$, exactly on schedule. And
across two full censuses, holes appeared only at lengths where the expected
count is small ($8$ and $16$), never at $32$ or beyond.

A first-moment heuristic therefore predicts no counterexample at any order
--- and the known escape from such heuristics is algebraic rigidity, the
loophole through which Moore graphs and cages exist. If the conjecture is
false, the witness is a highly structured object with spectral holes at
consecutive large powers of two. One general mechanism for building such
holes is a voltage lift certified through twisted non-backtracking spectra
(no identity-voltage non-backtracking closed $L$-walk in the base $\Rightarrow$
no $\Cyc L$ in the lift). The campaign's companion theoretical result
bounds this entire strategy:

\begin{theorem}[trace-hole Moore bound; proof in the repository notes]
\label{thm:tracemoore}
Let $G$ be a connected cubic multigraph on $n$ vertices, $H$ its Hashimoto
(non-backtracking) matrix, $L \ge 2$ even, $s = 2^{L/2}$. If
$\operatorname{tr}(H^L) = 0$ then
\[
n \;\ge\; \frac{(s+1)^2}{2s-1} \;\approx\; 2^{L/2-1},
\]
and twice that if $G$ is bipartite. In particular
$\operatorname{tr}(H^{16})=0$ forces $n \ge 130$,
$\operatorname{tr}(H^{32})=0$ forces $n \ge 32{,}770$, and
$\operatorname{tr}(H^{64})=0$ forces $n > 2.1\times10^9$.
\end{theorem}

Since $\operatorname{tr}(H^L) = 0$ implies $\Cyc d$-freeness for every
$d \mid L$, and a voltage lift is $\Cyc L$-free-by-walk-cancellation exactly
when its Hashimoto trace vanishes at $L$, walk-certified constructions
cannot reach $\Cyc{64}$ below two billion vertices. The bound is sharp in
flavor: the smallest cubic graph with $\operatorname{tr}(H^8)=0$ has exactly
$28$ vertices ($21$ witnesses, from a complete enumeration cross-checked
against Markstr\"om's counts), and the smallest known with
$\operatorname{tr}(H^{16})=0$ has $630$ (census record; bound says
$\ge 130$). Separately, an exact sweep of voltage lifts over every cubic
multigraph base on at most $6$ vertices (including loops and semi-edges) and
the $71$ loopless bases on $8$ vertices, against the stated group library ---
more than $2.3\times10^{10}$ configurations covering the window
$n \in [46,126]$ ---
found exactly three isomorphism classes avoiding
$\{\Cyc4,\Cyc8,\Cyc{16}\}$, each containing $\Cyc{32}$: \emph{no lift
counterexample exists below $128$ vertices in this scope}. In a different
companion direction, the conjecture's falsity is \emph{equivalent} to the
existence of an ``$\alpha$-block'' --- a $2$-connected graph with at most
one vertex of degree $2$ whose cycle spectrum avoids every power of two up
to its order --- a strictly more permissive search target; the exhaustive
ladder for these is empty through $31$ vertices
(\texttt{worktree-other-approaches} branch notes). C2030.1's cycle-only hole (walk count
$194{,}880$, cycle count $0$) shows a second, subtler mechanism exists
that the walk criterion cannot even see; the campaign obtained no general
classification of that mechanism.

\textbf{Unresolved quantities and scope limits.}
(1) The smallest cubic graph with no $\Cyc4,\Cyc8,\Cyc{16}$ remains in
$[50,82]$.
(2) The census exhibits both walk-cancellation and cycle-only holes at $16$;
no construction scaling these mechanisms to $32$ was obtained.
(3) The trace-hole bound ($n \ge 130$ for $L=16$) and the record
($n = 630$) leave a gap.
(4) The general minimum-order frontier established here is $36$.

\section{Data and reproducibility}
\label{sec:data}

The repository archives the source code, compact validation and result
records, the full campaign report with per-run wall/CPU times, and the graph
artifacts listed in the results directory. Bulk source censuses, most
per-shard solver logs, and long-cycle witness files are not duplicated in the
repository; the report records how those inputs and completion checks were
validated. The repository history timestamps each result, and the final
$n=48$ archive includes its monolith log, complete cube ledgers, aggregate
records, and drivers. An audited aggregate CPU total is not claimed.

\section*{Acknowledgments}

This campaign was carried out in collaboration with Claude (Anthropic),
which designed and operated the searches, wrote the toolchain, and drafted
this note under the author's direction. Source, graph artifacts, and compact
execution records are retained in the repository; the long-running
computations are reproducibility claims rather than formal proof
certificates. We thank Avery Carr, Arjun Balaji, and Klas
Markstr\"om, whose work this campaign builds on, and the maintainers of
\texttt{nauty}, SMS, CaDiCaL, and the Glasgow subgraph solver.

\begin{thebibliography}{99}

\bibitem{EG95} P.~Erd\H{o}s,
\emph{Some old and new problems in various branches of combinatorics},
Discrete Math.\ \textbf{165/166} (1997), 227--231. See also
\texttt{erdosproblems.com/64}.

\bibitem{DS01} D.~Daniel and S.~E.~Shauger,
\emph{A result on the \EG{} conjecture in planar graphs},
Congr.\ Numer.\ \textbf{153} (2001), 129--139.

\bibitem{Shauger98} S.~E.~Shauger,
\emph{Results on the \EG{} conjecture in $K_{1,m}$-free graphs},
Congr.\ Numer.\ \textbf{134} (1998), 101--105.

\bibitem{HK13} C.~C.~Heckman and R.~Krakovski,
\emph{\EG{} conjecture for cubic planar graphs},
Electron.\ J.\ Combin.\ \textbf{20}(2) (2013), \#P7.

\bibitem{P8free} \emph{\EG{} conjecture for $P_8$-free graphs},
Graphs Combin.\ \textbf{38} (2022). % Author list absent from the archived citation.

\bibitem{P10free} \emph{The \EG{} conjecture holds for $P_{10}$-free
graphs}, Discrete Math.\ (2024). % Author list absent from the archived citation.

\bibitem{LM23} H.~Liu and R.~Montgomery,
\emph{A solution to Erd\H{o}s and Hajnal's odd cycle problem},
J.\ Amer.\ Math.\ Soc.\ \textbf{36} (2023), 1191--1234.

\bibitem{Markstrom04} K.~Markstr\"om,
\emph{Extremal graphs for some problems on cycles in graphs},
Congr.\ Numer.\ \textbf{171} (2004), 179--192.

\bibitem{Carr26} A.~Carr,
\emph{Every minimal counterexample to the \EG{} conjecture is predominantly
cubic}, arXiv:2605.22844 (2026).

\bibitem{Balaji26} A.~Balaji,
\emph{erdos-gyarfas-min-degree-3} (software and SAT frontier $n \le 31$),
\url{https://github.com/ArjunBalaji79/erdos-gyarfas-min-degree-3} (2026).

\bibitem{nauty} B.~D.~McKay and A.~Piperno,
\emph{Practical graph isomorphism, II},
J.\ Symbolic Comput.\ \textbf{60} (2014), 94--112.

\bibitem{SMS} M.~Kirchweger and S.~Szeider,
\emph{SAT modulo symmetries for graph generation},
CP 2021; extended version J.\ Artif.\ Intell.\ Res.\ (2024).

\bibitem{Glasgow} C.~McCreesh, P.~Prosser, and J.~Trimble,
\emph{The Glasgow subgraph solver: using constraint programming to tackle
hard subgraph isomorphism problem variants}, ICGT 2020.

\bibitem{Conder} M.~Conder,
\emph{Trivalent (cubic) symmetric graphs}, census files,
\url{https://www.math.auckland.ac.nz/~conder/}.

\bibitem{PSV13} P.~Poto\v{c}nik, P.~Spiga, and G.~Verret,
\emph{Cubic vertex-transitive graphs on up to 1280 vertices},
J.\ Symbolic Comput.\ \textbf{50} (2013), 465--477.

\bibitem{SV08} B.~Sudakov and J.~Verstra\"ete,
\emph{Cycle lengths in sparse graphs},
Combinatorica \textbf{28} (2008), 357--372.

\end{thebibliography}

\end{document}
